La luz no viaja al azar. Desde hace siglos se sabe que cuando un rayo rebota en un espejo o cambia de medio, por ejemplo al pasar del aire al agua, lo hace siguiendo reglas muy precisas, y esas reglas se conocen como la ley de la reflexión y la ley de Snell.

Lo curioso es que ambas leyes salen de una sola idea.

Pierre de Fermat, un matemático francés del siglo XVII, propuso que la luz, para ir de un punto a otro, toma siempre el camino que le cuesta menos tiempo. A primera vista suena casi a que la luz pensara antes de moverse, pero en realidad es una forma muy elegante de resumir cómo se comporta, y además permite deducir las leyes de la óptica geométrica con un poco de cálculo.

En esta práctica no usamos láseres ni espejos reales. Todo se hizo dibujando. Con una regla y lápices de colores construimos varios caminos posibles entre dos puntos que pasaban por una línea de espejo, medimos cada uno y buscamos cuál era el más corto, para después comparar ese resultado con lo que predice la teoría.
